%% GRAFT — Grounding Region Annotations into Fine-tuned Transformers
%% NeurIPS 2024 style — compile with pdflatex
%%
%% \usepackage{neurips_2024}  -- download from:
%% https://neurips.cc/Conferences/2024/PaperInformation/StyleFiles
%%
%% If you don't have it, replace \documentclass line with:
%%   \documentclass[12pt,a4paper]{article}
%%   \usepackage[margin=1.5in]{geometry}
%% and remove the [preprint]{neurips_2024} package line.

\documentclass{article}
\PassOptionsToPackage{numbers,compress}{natbib}
\usepackage[preprint]{neurips_2024}

\usepackage[utf8]{inputenc}
\usepackage[T1]{fontenc}
\usepackage{hyperref}
\usepackage{url}
\usepackage{booktabs}
\usepackage{amsmath}
\usepackage{amssymb}
\usepackage{microtype}
\usepackage{graphicx}
\usepackage{multirow}
\usepackage{xcolor}
\usepackage{array}
\usepackage{caption}
\usepackage{float}
\usepackage{placeins}

\definecolor{darkgreen}{rgb}{0.0,0.45,0.0}
\definecolor{darkred}{rgb}{0.65,0.0,0.0}

\newcommand{\pg}{\textsc{pg}}
\newcommand{\ie}{\textit{i.e.}}
\newcommand{\eg}{\textit{e.g.}}

\title{GRAFT: Grounding Region Annotations into Fine-tuned Transformers}

\author{Siddharth Raj\\
  University of Chicago\\
  \href{https://github.com/Ani0202/region-grounded}{\texttt{sidraj@uchicago.edu}}
  \And
  Aniket Agrawal\\
  University of Chicago\\
  \href{https://github.com/Ani0202/region-grounded}{\texttt{agrawalaniket@uchicago.edu}}
}

\begin{document}
\maketitle

%% ─────────────────────────────────────────────────────────────────────────────
\begin{abstract}
Vision-language models trained on global image-caption contrastive objectives
produce patch features with little spatial structure: individual tokens know
what an image depicts, but not \emph{where}.  We ask whether a small amount of
region-level supervision can teach an explicit spatial grounding signal without
retraining the backbone. Concretely, we fine-tune SigLIP-B/16-384 using LoRA
(0.14\% of the parameters) with a combined global and FILIP-style region loss,
evaluated on the Flickr30k Pointing Game (PG).  Our main findings are: (1) the
dominant source of improvement is not the loss formulation or model capacity,
but \textbf{eliminating train/eval representation mismatches}---aligning the
image path (MaskCLIP bypass) and text path (EOS token) raised the 3-seed mean
from the baseline 14.74\% to 20.76\,$\pm$\,0.41\% with 6$\times$ lower
variance; (2) five epochs of training with per-epoch cosine restarts pushes a
single-seed result to \textbf{24.44\%} val PG / \textbf{26.16\%} test PG;
(3) auto-generated annotations (Florence-2) match human-drawn ones only when
the annotation vocabulary matches the evaluation distribution---switching from
verbose region captions to caption-grounded phrases closes an 8.25 pp gap and
reaches \textbf{23.35\%} PG.
\end{abstract}

%% ─────────────────────────────────────────────────────────────────────────────
\section{Introduction}
\label{sec:intro}

Modern vision-language models such as CLIP~\cite{radford2021clip} and
SigLIP~\cite{zhai2023siglip} learn to align images and text through contrastive
objectives applied at the image level. The internal patch tokens of these models
therefore encode \emph{what} is in the image but carry almost no information
about \emph{where} individual concepts are. This spatial blindness is a
fundamental limitation: models that cannot ground noun phrases to image regions
cannot support referring expression comprehension, dense captioning, or any
downstream task that requires spatial reasoning about language.

The standard remedy is to train or fine-tune on region-level annotations
(\eg, Flickr30k Entities~\cite{plummer2015flickr30k}, RefCOCO), which pair noun
phrases with bounding boxes. Full retraining is expensive; parameter-efficient
fine-tuning via LoRA~\cite{hu2022lora} is cheap. The question is how much
spatial grounding can be instilled in a globally-trained model by adding a
region loss to a LoRA fine-tuning run.

\textbf{Why SigLIP?} SigLIP replaces CLIP's softmax contrastive loss with
sigmoid contrastive loss, which trains on (image, text) pairs independently
rather than as a batch-level denominator.  This produces a stronger
image-text matching model, but the sigmoid objective has \emph{even less}
spatial bias than CLIP's softmax, because the denominator that might otherwise
encourage patch-level diversity is absent. SigLIP is a harder test than CLIP;
success here implies the approach generalises.

\textbf{What we build.} We call our approach GRAFT (Grounding Region
Annotations into Fine-tuned Transformers). Given a frozen SigLIP-B/16-384
backbone, we attach LoRA adaptors to the query and value projections of both
the vision and text encoders, and train with a loss that combines the original
global image-caption objective with a FILIP-style~\cite{yao2021filip}
phrase-to-patch alignment term. We evaluate with the Pointing Game
(PG)~\cite{zhang2018top}: for each annotated (phrase, bounding box) pair, the
argmax patch must fall inside the GT box.

\textbf{Core findings.}
\begin{enumerate}
\item \textbf{Representation mismatches matter more than loss design.} We
  identified and fixed two train/eval mismatches (image path, text path) that
  together account for the bulk of PG improvement. Ablations varying the loss
  weight, the LoRA target set, and the region loss formulation all returned to
  the same mean once the mismatches were corrected.
\item \textbf{Longer training with cosine restarts finds better basins.} Five
  epochs with per-epoch restarts exhibit a staircase dynamic: epochs 1--3
  reproduce the 1-epoch result; a step change occurs in epoch 4, lifting PG by
  3--5 pp.
\item \textbf{Annotation vocabulary is a critical design choice for auto-labels.}
  Florence-2's dense region captions (verbose descriptions) failed because the
  training vocabulary diverged from evaluation phrases. Switching to
  caption-grounded phrases recovered the full gain.
\end{enumerate}

%% ─────────────────────────────────────────────────────────────────────────────
\section{Background and Related Work}
\label{sec:background}

\paragraph{Spatial grounding from global VLMs.}
MaskCLIP~\cite{zhou2022maskclip} showed that replacing the last self-attention
layer of a CLIP vision encoder with a value-only projection (\ie, skipping
inter-patch mixing) yields patch features that are better structured for dense
prediction, without any training.  SCLIP~\cite{wang2023sclip} used
self-self attention to redistribute patch features based on spatial context.
Our work is complementary: we keep the MaskCLIP bypass as a fixed
representation choice and ask what supervised fine-tuning adds on top.

\paragraph{Region loss formulations.}
FILIP~\cite{yao2021filip} introduced token-level maximum similarity for
image-text alignment: instead of pooled features, take the maximum cosine
similarity between individual image and text tokens.  We adapt this to the
phrase-to-patch setting: for a phrase $p$ and image $I$, the region loss
requires $\max_{n \in \text{box}} \cos(f_n, e_p)$ to exceed
$\max_{n \notin \text{box}} \cos(f_n, e_p)$, where $f_n$ is the $n$-th patch
feature and $e_p$ is the phrase embedding.

\paragraph{Automatic annotation with Florence-2.}
Florence-2~\cite{xiao2024florence2} is a generalist vision-language model
supporting a wide range of annotation tasks through text prompts.  Its
\texttt{DENSE\_REGION\_CAPTION} task produces (bounding box, verbose
description) pairs; its \texttt{CAPTION\_TO\_PHRASE\_GROUNDING} task takes a
caption as input and grounds the noun phrases in it to bounding boxes.  We use
both, and show that vocabulary alignment between the annotation task and the
evaluation task is critical.

\paragraph{LoRA.}
LoRA~\cite{hu2022lora} freezes the pre-trained weight matrix $W_0 \in
\mathbb{R}^{d \times k}$ and adds a low-rank update $\Delta W = W_0 + B A$,
where $B \in \mathbb{R}^{d \times r}$, $A \in \mathbb{R}^{r \times k}$, and
$r \ll \min(d, k)$.  At inference, the delta is merged: $W_{\text{eff}} = W_0
+ \frac{\alpha}{r} B A$.  We use $r = 4$, $\alpha = 16$ (scale = 4), targets
$\{q\_\text{proj}, v\_\text{proj}\}$ in both encoders, yielding 295K trainable
parameters (0.14\% of the 200M-parameter backbone).

%% ─────────────────────────────────────────────────────────────────────────────
\section{Approach}
\label{sec:approach}

\subsection{Training Objective}
\label{sec:loss}

The total loss is:
\begin{equation}
\mathcal{L} = \mathcal{L}_{\text{global}} + \lambda \cdot \mathcal{L}_{\text{region}}
\label{eq:loss}
\end{equation}

\paragraph{Global loss $\mathcal{L}_{\text{global}}$.}
Standard SigLIP sigmoid contrastive loss between the MHAP pooler output of the
vision encoder and the EOS token of the text encoder, over image-caption pairs
in the batch.  The pooler output (\texttt{pooler\_output}), not
\texttt{mean(last\_hidden\_state)}, is critical: mean-pooling the patch tokens
directly in the loss collapses spatial diversity (Exp-01, Section~\ref{sec:exps}).

\paragraph{Region loss $\mathcal{L}_{\text{region}}$.}
For each (image, phrase, bounding box) triple $(I_j, p_i, B_{ij})$ in a region
batch, let $\hat{f}_n$ be the L2-normed patch feature at position $n$, and
$\hat{e}_i$ be the L2-normed EOS phrase embedding. The in-box and out-of-box
maximum similarities are:
\begin{equation}
s_{\text{in}} = \max_{n : n \in B_{ij}} \hat{f}_n \cdot \hat{e}_i,
\qquad
s_{\text{out}} = \max_{n : n \notin B_{ij}} \hat{f}_n \cdot \hat{e}_i
\end{equation}
The loss is sigmoid contrastive over an $N_{\text{region}} \times N_{\text{region}}$
matrix of $(s_{\text{in}}, s_{\text{out}})$ comparisons, matching the SigLIP
loss form. We set $\lambda = 0.5$; increasing it destabilises training
(Section~\ref{sec:ablations}).

\subsection{Representation Alignment}
\label{sec:alignment}

The most important design choices in GRAFT are not about the loss but about
ensuring the training and evaluation representations are identical.

\paragraph{Image path: MaskCLIP bypass during training.}
Evaluation reads patch features after the last self-attention layer is replaced
by a value-only projection:
\[
\text{bypass}(h) = W_O \, W_V \, h
\]
where $h$ is the input hidden state.  If training uses the full self-attention
forward pass but evaluation reads the bypass, the gradients sharpen features
under one representation while metrics measure another. We install the bypass
\emph{before training begins}, so the gradient flows through the same
$v\_\text{proj}$-only path that evaluation reads.  This single fix yields the
largest gain in the project (+7.34 pp on PG; Section~\ref{sec:img-fix}).

\paragraph{Text path: EOS token alignment.}
Early versions used \texttt{mean(last\_hidden\_state)} over all phrase tokens
in the region loss, while evaluation reads only the EOS position
(\texttt{last\_hidden\_state[:, -1, :]}). These are different vectors. We
replace the mean with the EOS token throughout training, creating a
byte-for-byte match with the evaluation text vector. This collapses
single-seed PG variance by $6\times$ and improves the 3-seed mean by +2.08 pp
(Section~\ref{sec:eos-fix}).

\subsection{Best-PG Snapshot}
\label{sec:snapshot}

The region loss can be satisfied by either a sharp argmax-friendly solution or
a broad coverage-friendly one; the optimizer drifts between these basins during
training. We address this with a lightweight best-checkpoint mechanism: during
training, we evaluate PG on the 200-image validation split every 200 steps,
snapshot the LoRA weights ($\sim$1 MB CPU memory) when PG improves, and
restore the best snapshot at the end of training. This costs zero extra compute
and provides free insurance against post-peak drift on every run.

\subsection{Datasets and Evaluation Protocol}
\label{sec:data}

\paragraph{Training data.} Flickr30k~\cite{young2014image} consists of 31,000
images with five captions each. We use the official train split (29,000 images)
for all training runs, optionally held out 10\% as a test set.  For the human
annotation condition (M\textsubscript{human}), bounding boxes come from
Flickr30k Entities~\cite{plummer2015flickr30k}, which provides human-drawn
boxes for noun phrases in the captions (average 5.7 phrases/image).  For
automatic annotation conditions, we use Florence-2 to generate (phrase, box)
pairs from the same images.

\paragraph{Evaluation.} All models are evaluated on a fixed 200-image subset
of the Flickr30k validation split (seed=42), containing 2,124 phrase-box pairs
(10.6 pairs/image on average). The \textbf{Pointing Game (PG)} score is the
fraction of pairs where the centre of the highest-similarity patch falls inside
the GT bounding box. We also measure \textbf{Recall@1} (the fraction of pairs
where the tight box around the top-$k$ selected patches overlaps the GT box at
IoU $\geq 0.5$), using $k \in \{10, 25, \text{halfmax}\}$.  A held-out test
set of 2,900 images (28,467 phrase-box pairs, SE $\approx$ 0.24 pp) provides
statistically reliable generalisation estimates.

\paragraph{Frozen baseline (B0).} SigLIP-B/16-384 with MaskCLIP bypass at
evaluation, no fine-tuning. B0 PG = \textbf{14.74\%}. The gap between PG
(14.74\%) and Recall@1 (6.83--7.91\%) indicates the model has a weak
directional signal but no coherent spatial clustering: the argmax occasionally
lands near the right region, but the surrounding high-scoring patches scatter
across the image.

%% ─────────────────────────────────────────────────────────────────────────────
\section{Experiments}
\label{sec:exps}

We ran 15 experiments progressing from a broken initial recipe to the final
locked configuration. We describe them grouped by the mechanistic insight each
contributed.

\subsection{Initial Failures: MHAP Pooler Fix (Exp-01, 02)}
\label{sec:initial}

The first training run used mean-pooling of all patch tokens for the global
loss. Both the model with (M\textsubscript{human}) and without (B1) the region
loss regressed below the frozen baseline on VOC mIoU, and PG also regressed.
Root cause: mean-pooling pushes all 576 patch features uniformly toward the
caption, collapsing the spatial diversity that grounding requires.

Switching to the MHAP \texttt{pooler\_output} (a dedicated max-pooling head in
SigLIP) fixed the collapse: patch features were freed to be spatially diverse.
However, PG still sat at 13.42\% (below B0) after this fix alone, for a
different reason traced in the next experiment.

\subsection{Image-Path Alignment: MaskCLIP Bypass During Training (Exp-03)}
\label{sec:img-fix}

With the MHAP fix in place but without the MaskCLIP bypass during training, PG
= 13.42\% (below B0) while Recall@1 top-10 improved slightly (+0.28 pp). This
mixed result is the signature of a \emph{train/eval path mismatch}: training
read patch features through the full last self-attention layer, but evaluation
read them through the bypass. The gradient sharpened patches under one
representation; the argmax metric measured the other. The Recall@1 improvement
confirmed the model was learning something real---it was just invisible to
argmax localization.

Installing the MaskCLIP bypass before training produced a step change:
\textbf{PG = 22.08\% (+7.34 pp over B0)} on the first run. The bypass fix also
changed the \emph{shape} of the response: training through full attention
produced broad, weak activation across the bounding box (good for Recall@1
coverage, bad for argmax); training through the bypass produced a single sharp
spike on one patch inside the box---the behaviour the max-pool region loss
actually rewards.

\subsection{Variance Quantification (Exp-06)}
\label{sec:variance}

The 22.08\% result was a single run. Subsequent runs with ``identical''
hyperparameters but different DataLoader shuffle orders produced 18.41\% and
16.90\%---a 5.18 pp spread ($\approx 6\sigma$ on the binomial standard error).
We added full seed control (Python, NumPy, PyTorch CPU/CUDA, DataLoader
generator) and ran 3 seeds. Results:

\begin{center}
\begin{tabular}{lccccc}
\toprule
 & B0 & seed 0 & seed 1 & seed 2 & mean $\pm$ std \\
\midrule
PG & 14.74\% & 19.26\% & 15.87\% & 20.90\% & \textbf{18.68 $\pm$ 2.57\%} \\
R@1 halfmax & 7.58\% & 7.30\% & 7.30\% & 7.30\% & 7.30 $\pm$ 0.00\% \\
\bottomrule
\end{tabular}
\end{center}

The 22.08\% result was a $+1.3\sigma$ draw; the true mean improvement is
+3.94 pp. Crucially, R@1 halfmax is \emph{identical across all three seeds}
(to four decimal places), while PG varies by 5 pp. The recipe is learning a
stable broader grounding signal---the variance is entirely in which single
patch wins the argmax for borderline phrases. From this point, we adopted a
``leader-seed'' protocol: test new ideas on seed=2 first; run all 3 seeds only
if seed=2 exceeds 21.25\% (mean + $\sigma$).

\subsection{Ablations: Loss Weight, Capacity, and Aggregation (Exp-07--10)}
\label{sec:ablations}

We tested three hypotheses for pushing PG above the 18.68\% mean, all on
seed=2 against the 20.90\% leader baseline.

\textbf{$\lambda$ 0.5 $\to$ 1.0 (Exp-07).} PG fell to 17.18\% (--3.72 pp).
Doubling the region loss weight broadened the patch response rather than
sharpening it: more in-bbox patches at medium similarity satisfies the max-pool
objective just as well as a sharp spike, and the optimizer took the broader
path. The global loss, whose weight halved in relative terms, provides
structural regularisation that constrains patch features to a coherent
manifold; removing it caused a mid-training collapse (PG below B0 at step 400).
Notably, R@1 top-25 \emph{increased} to 7.58\% (+0.56 pp over v5 seed=2)
even as PG fell---confirming the broadening interpretation: more in-bbox
patches at medium similarity, but not one strong argmax spike.

\textbf{k-proj LoRA, q+v $\to$ q+k+v (Exp-08).} Adding key projections to the
LoRA target set (+50\% trainable params, 442K total) produced the highest
checkpoint of any run: 22.69\% PG at step 400---but the model then drifted to
17.80\% by end of epoch. More capacity let the model find a sharp grounding
pattern faster, but also gave the optimizer more freedom to walk away from it.
The \emph{peak} improved; the \emph{final result} did not.

\textbf{Top-K mean region loss (Exp-09, 10).} Adding the best-PG snapshot to
the q+k+v run (Exp-09) confirmed that the snapshot mechanism works: final PG
= 19.96\% (saved from the 17.51\% end-of-epoch state). Replacing the max-pool
region loss with a top-3 mean (Exp-10) produced PG = 19.35\%. Neither move
changed the 3-seed mean. We concluded that the bottleneck was not on the image
or loss side. As shown in Table~\ref{tab:ablations}, R@1 halfmax is 7.30\% on
every single run (identical to 4 decimal places)---all four ablations learn the
same broad grounding signal; the variance is entirely in PG (argmax).

\begin{table}[H]
\centering
\caption{Ablation results on seed=2. All runs use the v5 base recipe
(q+v LoRA, $\lambda{=}0.5$, 1 ep, best-PG snapshot from Exp-09 onward)
with one change each. R@1 halfmax is constant across all variants.}
\label{tab:ablations}
\setlength{\tabcolsep}{5pt}
\begin{tabular}{lcccc}
\toprule
Variant & PG & R@1 top-10 & R@1 top-25 & R@1 halfmax \\
\midrule
v5 seed=2 (baseline)      & 20.90\% & 5.98\% & 7.02\% & 7.30\% \\
v6: $\lambda{=}1.0$        & 17.18\% & 5.93\% & \textbf{7.58\%} & 7.30\% \\
v7: +k-proj LoRA           & 17.80\% & 5.46\% & 7.25\% & 7.30\% \\
v8: +best-PG snap          & 19.96\% & ---    & ---    & 7.30\% \\
v9: top-3 mean loss        & 19.35\% & 5.50\% & 7.18\% & 7.30\% \\
\bottomrule
\end{tabular}
\end{table}

\subsection{Text-Path Alignment: EOS Token Fix (Exp-11)}
\label{sec:eos-fix}

The image-side mismatch (training through full attention vs. MaskCLIP bypass)
had been fixed in Exp-03. An analogous mismatch on the text side had persisted
throughout:
\begin{itemize}
\item \textbf{Training}: region loss used \texttt{mean(last\_hidden\_state)}
  over all phrase tokens.
\item \textbf{Evaluation}: PG used \texttt{last\_hidden\_state[:, -1, :]}---the
  EOS position only.
\end{itemize}
The region loss was teaching patches to align with the mean-over-tokens
representation, but PG measured alignment with the EOS representation. Fixing
this is mechanistically identical to the image-side fix, applied to the text
encoder.

We replaced \texttt{mean-over-tokens} with the EOS token throughout training,
reverted to q+v LoRA (k-proj correlated with drift), and kept the best-PG
snapshot. Results across 3 seeds:

\begin{center}
\begin{tabular}{llcccc}
\toprule
 & & seed 0 & seed 1 & seed 2 & mean $\pm$ std \\
\midrule
\multirow{2}{*}{v5} & PG       & 19.26\% & 15.87\% & 20.90\% & 18.68 $\pm$ 2.57\% \\
                    & R@1 hmax & 7.30\%  & 7.30\%  & 7.30\%  & 7.30 $\pm$ 0.00\% \\
\midrule
\multirow{2}{*}{v10 (EOS fix)} & PG       & 20.57\% & 20.48\% & 21.23\% & \textbf{20.76 $\pm$ 0.41\%} \\
                               & R@1 hmax & 7.44\%  & 7.34\%  & ---     & --- \\
\bottomrule
\end{tabular}
\end{center}

The EOS fix raised the mean by +2.08 pp and collapsed the variance by $6\times$
(from 2.57 to 0.41 pp). Every seed now falls in the range [20.48, 21.23]---a
0.75 pp spread vs. v5's 5.03 pp spread. The best-PG snapshot eliminated
late-training drift entirely on every seed: peak PG and final PG are identical
for all three. This recipe (v10) is the locked M\textsubscript{human}
configuration.

\textbf{Mechanistic reading.} The EOS fix raised the \emph{peak}; the best-PG
snapshot locked it. Both ingredients contribute orthogonally. The drift itself
(still visible in raw trajectories) is a generic overfit dynamic in LoRA on
small data; the snapshot is the correct mitigation regardless of its cause.

\subsection{Generalisation Test (Exp-12)}
\label{sec:gentest}

The best-PG snapshot selects the checkpoint with the highest validation PG,
raising the question of whether the reported number is optimistically biased by
the selection step. We held out 10\% of the Flickr30k training split (2,900
images, 28,467 phrase-box pairs) as an independent test set, never used for
any training-time decision.

\begin{center}
\begin{tabular}{lcccccc}
\toprule
seed & val PG & test PG & $\Delta$ & R@1 top-10 & R@1 top-25 & R@1 halfmax \\
\midrule
0 & 19.92\% & 19.84\% & $-$0.08 & 5.74\% & 7.72\% & 7.25\% \\
1 & 21.61\% & 21.83\% & $+$0.22 & 6.59\% & 7.58\% & 7.34\% \\
2 & 17.94\% & 18.52\% & $+$0.58 & 6.54\% & 7.34\% & 7.25\% \\
\midrule
mean & 19.82$\pm$1.84\% & 20.07$\pm$1.67\% & $+$0.25 & --- & --- & 7.28\% \\
\bottomrule
\end{tabular}
\end{center}

The val$\to$test delta is effectively zero. Test PG is \emph{higher} than val
PG on average, and every seed's test PG is within 0.6 pp of its val PG. The
best-PG snapshot selects genuinely good checkpoints, not ones that happened to
score well on the validation split by chance. With 28,467 pairs (SE $\approx$
0.24 pp), the +5.33 pp test gain over B0 is a $\sim$22$\sigma$ signal.

\subsection{Five-Epoch Training (Exp-13)}
\label{sec:5ep}

With the locked v10 recipe, training plateaus or drifts after step 400--800
($\sim$half the epoch). We trained for 5 epochs with per-epoch cosine restarts
(\texttt{restarts=5}, each epoch its own warmup$\to$cosine$\to$0 cycle).

The PG trajectory for seed=1 (Table~\ref{tab:v11traj}) shows a clear
staircase pattern: epochs 1--3 reproduce the 1-epoch v10 result ($\sim$19--20\%
val PG); a step change occurs in epoch 4 ($+$3--5 pp), and the gain persists
through epoch 5. The per-epoch restarts kick the optimizer out of the epoch-1
basin into a deeper minimum.

\begin{table}[H]
\centering
\caption{Selected steps from the 5-epoch training trajectory (seed=1).}
\label{tab:v11traj}
\begin{tabular}{rlrr}
\toprule
Step & Epoch & Val PG & Test PG \\
\midrule
800  & 1 & 19.68\% & 19.92\% \\
1600 & 2 & 19.40\% & 20.74\% \\
2400 & 3 & 19.21\% & 21.41\% \\
3000 & 4 & 23.59\% & 24.59\% \\
\textbf{3400} & \textbf{5} & \textbf{24.44\%} & \textbf{25.26\%} \\
4000 & 5 & 24.11\% & 26.03\% \\
\bottomrule
\end{tabular}
\end{table}

The best-PG snapshot was taken at step 3400 (24.44\% val PG). The final test
PG on the full held-out set is \textbf{26.16\%} (7,447/28,467 pairs).
Table~\ref{tab:v11final} shows the full metric breakdown at the best-PG
checkpoint.

\begin{table}[H]
\centering
\caption{M\textsubscript{human} v11 final metrics at best-PG checkpoint
(step 3400, seed=1). R@1 top-25 = 8.29\% is the only trained model to
beat B0's 7.91\%, indicating genuine spatial coherence after 5 epochs.}
\label{tab:v11final}
\begin{tabular}{lrr}
\toprule
Metric & Val & Test \\
\midrule
Pointing Game     & \textbf{24.44\%} & \textbf{26.16\%} \\
R@1 top-10        & 5.93\%  & ---    \\
R@1 top-25        & \textbf{8.29\%}  & 7.67\% \\
R@1 halfmax       & 7.30\%  & 7.35\% \\
\midrule
\textit{B0 baseline} & & \\
~~Pointing Game   & 14.74\% & ---    \\
~~R@1 top-25      & 7.91\%  & ---    \\
~~R@1 halfmax     & 7.58\%  & ---    \\
\bottomrule
\end{tabular}
\end{table}

\subsection{Automatic Annotations: M\textsubscript{auto} and M\textsubscript{cpg} (Exp-14, 15)}
\label{sec:mauto}

\paragraph{M\textsubscript{auto}: Dense Region Captions (Exp-14).}
We replaced Flickr30k Entities with Florence-2 \texttt{DENSE\_REGION\_CAPTION}
annotations: automatically generated (bounding box, verbose description) pairs
(29,000 images, $\sim$9.8 regions/image). The locked v10 recipe was used
unchanged. Result: best-PG $\approx$15.1\% (barely above B0 = 14.74\%). PG
collapsed to $\sim$11.5\% at step 600 and recovered only partially.

\textbf{Diagnosis: train/eval vocabulary mismatch.} Florence-2
\texttt{DENSE\_REGION\_CAPTION} generates verbose descriptions
(\eg, ``a woman in a blue top standing near a fence''). Evaluation uses short
Flickr30k Entities phrases (\eg, ``woman'', ``the man''). SigLIP's
pre-trained representations can bridge this gap at initialisation ($\sim$15\%);
as region loss gradients accumulate, patch features align with the verbose
training vocabulary and diverge from the short-phrase evaluation distribution.
Bounding box \emph{geometry} from Florence-2 is accurate; the \emph{labels}
are the problem.

\paragraph{M\textsubscript{cpg}: Caption-to-Phrase Grounding (Exp-15).}
Florence-2 \texttt{CAPTION\_TO\_PHRASE\_GROUNDING} takes the Flickr30k caption
as input and returns noun phrases grounded to bounding boxes---the same caption
that drives the global loss, and noun phrases from the same distribution as
Flickr30k Entities evaluation phrases. We ran 5 epochs on seed=1.

\begin{center}
\begin{tabular}{lcc}
\toprule
Epoch & Val PG & Test PG \\
\midrule
1 & 22.60\% & 22.50\% \\
\textbf{2} & \textbf{23.35\%} & 22.79\% \\
3 & 21.99\% & 21.96\% \\
4 & 22.65\% & \textbf{23.15\%} \\
5 & 22.03\% & 22.13\% \\
\bottomrule
\end{tabular}
\end{center}

Best val PG = 23.35\% (epoch 2); best test PG = 23.15\% (epoch 4).
The +8.25 pp gain over M\textsubscript{auto} (same Florence-2 model, same
images, same recipe, only labels changed) isolates annotation vocabulary as the
sole cause of M\textsubscript{auto}'s failure.

M\textsubscript{cpg} (23.35\%) also outperforms M\textsubscript{human} v10
(20.48\%, seed=1) by +2.87 pp. Florence-2 grounds phrases from all five
Flickr30k captions per image, producing denser supervision than the human
Entities annotations, which cover a fixed set of phrases per image. When
vocabulary is controlled, automatic annotations at scale can exceed human-drawn
ones. Table~\ref{tab:main} summarises all models; the full experiment log is in
Appendix~\ref{app:explog}.

\begin{table}[H]
\centering
\caption{Main results on the Flickr30k Pointing Game. All models use
SigLIP-B/16-384 with MaskCLIP bypass at evaluation. Val: 200 images, 2,124
pairs. Test: 2,900 held-out images, 28,467 pairs. $^\dagger$Single seed;
mean $\pm$ std pending. R@1 = Recall@1 at IoU$\geq$0.5; halfmax = all
patches with similarity $\geq 0.5 \times$ max; ``---'' = not measured.}
\label{tab:main}
\setlength{\tabcolsep}{4pt}
\begin{tabular}{llllll}
\toprule
Model & Config & Val PG & Test PG & R@1 top-25 & R@1 halfmax \\
\midrule
B0 (frozen) & — & 14.74\% & --- & 7.91\% & 7.58\% \\
\midrule
M\textsubscript{human} v5  & FILIP, 1 ep        & 18.68$\pm$2.57\% & ---               & 7.12$\pm$0.12\% & \textbf{7.30$\pm$0.00\%} \\
M\textsubscript{human} v10 & EOS+snap, 1 ep     & \textbf{20.76$\pm$0.41\%} & 20.07$\pm$1.67\% & ---      & 7.28\%$^\ddagger$ \\
M\textsubscript{human} v11 & EOS+snap, 5 ep     & 24.44\%$^\dagger$ & 26.16\%$^\dagger$ & \textbf{8.29\%} & 7.30\% \\
\midrule
M\textsubscript{auto}      & Dense cap., 1 ep   & $\sim$15.1\%$^\dagger$ & ---          & ---             & --- \\
M\textsubscript{cpg}       & Phrase ground., 5 ep & 23.35\%$^\dagger$ & 23.15\%$^\dagger$ & ---           & --- \\
\bottomrule
\end{tabular}
{\small $^\ddagger$Mean of seeds 0 and 1 only (seed 2 not recorded).}
\end{table}

%% ─────────────────────────────────────────────────────────────────────────────
\section{Results and Discussion}
\label{sec:results}

\paragraph{The two biggest gains come from alignment, not architecture.}
The image-path fix (v1$\to$v2, +7.34 pp on one seed, +3.94 pp on 3-seed mean)
and the text-path fix (v5$\to$v10, +2.08 pp mean, $6\times$ variance
reduction) together account for the majority of total improvement. All ablation
experiments varying the loss formulation, loss weight, and LoRA target set
returned to the same mean ($\sim$18--21\%) once the mismatches were in place.
This suggests that the bottleneck in adapting globally-trained VLMs for spatial
grounding is not loss design but representation consistency.

\paragraph{Variance is concentrated in argmax, not coverage.}
The R@1 halfmax metric (all patches above half the maximum similarity) is
\emph{identical} across seeds (7.30 $\pm$ 0.00\% on v5), while PG varies by
2.57 pp. The model learns a stable broad grounding signal; the 3-5 pp PG
variance is purely in which single patch wins the argmax for borderline
phrases. R@1 halfmax is the more reliable metric for monitoring whether the
recipe is learning anything at all; PG is the more informative metric for
comparing strong models.

\paragraph{Longer training unlocks a second optimisation basin.}
The staircase pattern in v11 (Table~\ref{tab:v11traj}) suggests that the
1-epoch recipe converges to a local minimum from which small perturbations
(changing $\lambda$, adding k-proj LoRA) cannot escape. Per-epoch cosine
restarts provide the momentum to leave that basin at epoch 4. This is a
practical insight: when ablations around a recipe all fail, more training may
be more productive than further hyperparameter search.

\paragraph{Vocabulary alignment dominates annotation quality for auto-labels.}
The 8.25 pp gap between M\textsubscript{auto} and M\textsubscript{cpg} is
entirely attributable to label vocabulary. Florence-2 can produce accurate
bounding boxes; the constraint is that the phrase labels must come from the same
distribution as evaluation phrases. This is a non-trivial requirement when
evaluation uses a fixed dataset with its own vocabulary (Flickr30k Entities):
any auto-labelling pipeline must either match that vocabulary or rely on the
pre-trained model's cross-vocabulary generalisation not degrading under
fine-tuning.

%% ─────────────────────────────────────────────────────────────────────────────
\section{Qualitative Results}
\label{sec:qualitative}

Figure~\ref{fig:viz} shows the most dramatic example from our validation scan:
a case where both fine-tuned models clearly outperform the frozen baseline.
Each column shows the plasma-normalised cosine-similarity heatmap blended over
the image, the ground-truth bounding box (white rectangle), and the argmax
patch centre ($\bigstar$; {\color{green!60!black}green} = hit,
{\color{red}red} = miss). All six cherry-picked images are in
Appendix~\ref{app:viz} (Figures~\ref{fig:app1}--\ref{fig:app6}).

\begin{figure}[p]
\centering
\includegraphics[width=\linewidth,height=0.82\textheight,keepaspectratio]{pg_viz_04_1507563902}
\caption{\textbf{Most dramatic case.} Image \texttt{1507563902.jpg}
($n{=}3$ phrases).
B0: 33\%\enspace M\textsubscript{human}: \textbf{100\%}\enspace
M\textsubscript{auto}: \textbf{67\%}.
B0 localises only 1 of 3 phrases; M\textsubscript{human} hits all 3.
Columns: Frozen B0 (left) / M\textsubscript{human} (centre) /
M\textsubscript{auto} (right).
Plasma heatmap = patch-phrase cosine similarity; white box = GT bounding box;
$\bigstar$ = argmax patch ({\color{green!60!black}green} hit /
{\color{red}red} miss).}
\label{fig:viz}
\end{figure}

%% ─────────────────────────────────────────────────────────────────────────────
\section{Conclusion}
\label{sec:conclusion}

We have shown that LoRA fine-tuning with a region loss can teach meaningful
spatial grounding to a globally-trained VLM (SigLIP-B/16-384), raising the
Flickr30k Pointing Game from 14.74\% to 24.44\% val / 26.16\% test (single
seed, 5 epochs) and 20.76\,$\pm$\,0.41\% (3-seed mean, 1 epoch).  The key
lessons are:

\begin{enumerate}
\item \textbf{Fix mismatches first.} The two biggest gains in the project came
  from aligning the image path (MaskCLIP bypass during training) and text path
  (EOS token throughout) with evaluation. Neither required any architectural
  change or new data. Loss and architecture ablations had small or negative
  effect on the mean.
\item \textbf{Measure variance before drawing conclusions.} The first
  single-point result (22.08\%) was 1.3$\sigma$ above the true 3-seed mean.
  A proper variance study should precede any recipe change comparison.
\item \textbf{The global loss is structural, not optional.} Reducing its weight
  removes a constraint that keeps patch features in a coherent manifold;
  training destabilises.
\item \textbf{Best-PG snapshot is free insurance.} It costs $\sim$1 MB of CPU
  memory, zero compute, and eliminates end-of-training drift on every run.
\item \textbf{For auto-labels, vocabulary controls everything.} Florence-2
  bounding boxes are geometrically accurate; the critical design choice is
  which annotation task supplies the phrase labels.
\end{enumerate}

The strongest remaining open question is whether the 5-epoch result (24.44\%)
is robust across seeds---the $\sigma \approx 2.57$ pp observed in 1-epoch
training suggests meaningful variance may persist. Running seeds 0 and 2 for
both v11 and M\textsubscript{cpg} is the immediate next step.

\paragraph{Code availability.}
All training scripts, evaluation code, and experiment notebooks are available at
\url{https://github.com/Ani0202/region-grounded} (branch: \texttt{main}).

%% ─────────────────────────────────────────────────────────────────────────────
\section{Limitations and Future Work}
\label{sec:future}

\paragraph{Broader evaluation benchmarks.}
The Pointing Game on Flickr30k is a useful proxy but a narrow one: binary
hit/miss on 200 images with a fixed phrase vocabulary. RefCOCO and RefCOCO+
are the standard benchmarks for referring expression comprehension, covering a
different annotation style (relational expressions such as ``the person on the
left'') and a different image domain (MS-COCO). Evaluating GRAFT on these
would test whether the grounding signal transfers beyond the distribution it
was trained on.

\paragraph{Dense segmentation.}
VOC mIoU was tracked throughout as a guardrail and consistently regressed
(B0 = 1.43\%; all trained models stayed below). This was expected: the
FILIP max-pool loss rewards one sharp patch per phrase, not a spatially
coherent mask. Now that PG has reached 24.44\%, it is worth revisiting
whether the stronger grounding signal incidentally improves dense prediction,
or whether PG and segmentation mIoU are genuinely measuring orthogonal
properties of patch features.

\paragraph{Grounding at multiple granularities.}
All experiments used short noun phrases (\eg, ``the woman'', ``a red jacket'')
as the grounding target. An open question is whether the same recipe extends
to longer, relational descriptions (\eg, ``the person on the left wearing a
red jacket'') that require integrating evidence across multiple spatial
locations. Such descriptions may require a richer text representation than
a single EOS token and a loss formulation that rewards patch \emph{sets}
rather than individual patches.

%% ─────────────────────────────────────────────────────────────────────────────
\section*{References}
\begingroup
\renewcommand{\section}[2]{}
\begin{thebibliography}{99}

\bibitem{radford2021clip}
A.\ Radford, J.\ W.\ Kim, C.\ Hallacy, A.\ Ramesh, G.\ Goh, S.\ Agarwal,
G.\ Sastry, A.\ Askell, P.\ Mishkin, J.\ Clark, G.\ Krueger, and I.\ Sutskever.
Learning transferable visual models from natural language supervision.
In \emph{ICML}, 2021.

\bibitem{zhai2023siglip}
X.\ Zhai, B.\ Mustafa, A.\ Kolesnikov, and L.\ Beyer.
Sigmoid loss for language image pre-training.
In \emph{ICCV}, 2023.

\bibitem{hu2022lora}
E.\ J.\ Hu, Y.\ Shen, P.\ Wallis, Z.\ Allen-Zhu, Y.\ Li, S.\ Wang, L.\ Wang,
and W.\ Chen.
LoRA: Low-rank adaptation of large language models.
In \emph{ICLR}, 2022.

\bibitem{yao2021filip}
L.\ Yao, R.\ Huang, L.\ Hou, G.\ Lu, M.\ Niu, H.\ Xu, X.\ Liang, Z.\ Li,
X.\ Jiang, and C.\ Xu.
FILIP: Fine-grained interactive language-image pre-training.
In \emph{ICLR}, 2022.

\bibitem{zhou2022maskclip}
C.\ Zhou, C.\ C.\ Loy, and B.\ Dai.
Extract free dense labels from CLIP.
In \emph{ECCV}, 2022.

\bibitem{wang2023sclip}
F.\ Wang, J.\ Cui, Y.\ Kim, A.\ Guo, Y.\ Lu, and A.\ Yuille.
SCLIP: Rethinking self-attention for dense vision-language inference.
\emph{arXiv:2312.01597}, 2023.

\bibitem{plummer2015flickr30k}
B.\ A.\ Plummer, L.\ Wang, C.\ M.\ Cervantes, J.\ C.\ Caicedo,
J.\ Hockenmaier, and S.\ Lazebnik.
Flickr30k entities: Collecting region-to-phrase correspondences for richer
image-to-sentence models.
In \emph{ICCV}, 2015.

\bibitem{young2014image}
P.\ Young, A.\ Lai, M.\ Hodosh, and J.\ Hockenmaier.
From image descriptions to visual denotations: New similarity metrics for
semantic inference over event descriptions.
\emph{TACL}, 2014.

\bibitem{xiao2024florence2}
B.\ Xiao, H.\ Wu, W.\ Xu, X.\ Dai, H.\ Hu, Y.\ Lu, M.\ Zeng, C.\ Liu, and
L.\ Yuan.
Florence-2: Advancing a unified representation for a variety of vision tasks.
In \emph{CVPR}, 2024.

\bibitem{zhang2018top}
J.\ Zhang, Z.\ Lin, J.\ Brandt, X.\ Shen, and S.\ Sclaroff.
Top-down neural attention by excitation backprop.
\emph{IJCV}, 2018.

\end{thebibliography}
\endgroup

%% ─────────────────────────────────────────────────────────────────────────────
\FloatBarrier
\appendix

\section{Full Experiment Log}
\label{app:explog}

Table~\ref{tab:allexps} reports the complete sequence of experiments with their
primary PG result. Entries marked $\dagger$ are single seeds (seed=2 unless
noted); entries with $n=3$ report mean $\pm$ std.

\begin{table}[H]
\centering
\caption{All 15 experiments. $\lambda=0.5$ throughout unless noted. Best-PG
snapshot used from Exp-09 onward.}
\label{tab:allexps}
\setlength{\tabcolsep}{4pt}
\begin{tabular}{clll}
\toprule
Exp & Key change & Val PG & Notes \\
\midrule
01 & Baseline LoRA + mean-pool global & Regression & VOC mIoU regressed \\
02 & MHAP pooler fix & 13.42\%$^\dagger$ & PG below B0 \\
03 & MaskCLIP bypass during training & 22.08\%$^\dagger$ & +7.34 pp; path mismatch fixed \\
04 & 2 ep single cosine & $\leq$19.40\%$^\dagger$ & LR too high too long \\
05 & 2 ep cosine restart & 16.90--18.41\%$^\dagger$ & High variance identified \\
06 & Seed control, 3 seeds & 18.68 $\pm$ 2.57\% & True mean +3.94 pp \\
07 & $\lambda$: 0.5$\to$1.0 & 17.18\%$^\dagger$ & Broadens response; collapse \\
08 & +k-proj LoRA & 17.80\%$^\dagger$ (22.69\% peak) & Peak good, drifts \\
09 & +best-PG snapshot & 19.96\%$^\dagger$ & Snapshot works; k-proj neutral \\
10 & Top-K mean region loss & 19.35\%$^\dagger$ & No improvement \\
11 & EOS text alignment, 3 seeds & \textbf{20.76 $\pm$ 0.41\%} & $6\times$ variance reduction \\
12 & Generalisation test (90-10) & 19.82 $\pm$ 1.84\% & Test: 20.07 $\pm$ 1.67\%; no leakage \\
13 & 5 ep cosine restart & \textbf{24.44\%}$^\dagger$ & Test 26.16\%; staircase \\
14 & M\textsubscript{auto} dense caption & $\sim$15.1\%$^\dagger$ & Vocab mismatch failure \\
15 & M\textsubscript{cpg} phrase ground. & \textbf{23.35\%}$^\dagger$ & +8.25 pp vs.\ M\textsubscript{auto} \\
\bottomrule
\end{tabular}
\end{table}

\section{Ablation Detail: v5 Trajectory by Seed}
\label{app:v5traj}

\begin{table}[H]
\centering
\caption{PG at each evaluation step during 1-epoch training (v5). Binomial SE
$\approx 0.87$ pp at $p \approx 0.19$, $n=2124$.}
\begin{tabular}{rrrrr}
\toprule
Step & Epoch & seed 0 & seed 1 & seed 2 \\
\midrule
0   & 0 & 14.74 & 14.74 & 14.74 \\
200 & 1 & 16.43 & 16.15 & 16.48 \\
400 & 1 & 16.71 & 19.02 & 19.92 \\
600 & 1 & 19.30 & 14.92 & 20.62 \\
800 & 1 & 19.16 & 16.10 & 20.95 \\
end & 1 & 19.26 & 15.87 & 20.90 \\
\bottomrule
\end{tabular}
\end{table}

Seed 1's collapse to 14.92\% at step 600 (below B0) and partial recovery to
15.87\% is a consistent feature of the loss/metric landscape, not a randomness
artefact. It also occurred in Exp-05 (v4 attempt 2) and reflects a training
basin where the FILIP objective is satisfied by attention patterns that hurt
argmax localisation.

\section{M\textsubscript{human} v11 Full Trajectory}
\label{app:v11traj}

\begin{table}[H]
\centering
\caption{Full 5-epoch training trajectory for M\textsubscript{human} v11,
seed=1. Best-PG snapshot taken at step 3400.}
\begin{tabular}{rlrr}
\toprule
Step & Epoch & Val PG (\%) & Test PG (\%) \\
\midrule
200  & 1 & 16.34 & 15.04 \\
400  & 1 & 19.59 & 18.48 \\
600  & 1 & 19.63 & 19.92 \\
800  & 1 & 19.68 & 19.92 \\
1000 & 2 & 17.33 & 18.94 \\
1200 & 2 & 18.55 & 20.53 \\
1400 & 2 & 18.79 & 20.12 \\
1600 & 2 & 19.40 & 20.74 \\
1800 & 3 & 17.94 & 18.69 \\
2000 & 3 & 17.47 & 19.71 \\
2200 & 3 & 18.64 & 22.23 \\
2400 & 3 & 19.21 & 21.41 \\
2600 & 4 & 22.74 & 23.20 \\
2800 & 4 & 22.69 & 22.43 \\
3000 & 4 & 23.59 & 24.59 \\
3200 & 4 & 23.26 & 25.00 \\
\textbf{3400} & \textbf{5} & \textbf{24.44} & \textbf{25.26} \\
3600 & 5 & 23.26 & 24.79 \\
3800 & 5 & 24.15 & 25.72 \\
4000 & 5 & 24.11 & 26.03 \\
\bottomrule
\end{tabular}
\end{table}

\section{Locked Recipe}
\label{app:recipe}

\begin{table}[H]
\centering
\caption{GRAFT locked recipe (M\textsubscript{human} v10 / v11).}
\begin{tabular}{ll}
\toprule
Component & Setting \\
\midrule
Backbone & SigLIP-B/16-384, all parameters frozen \\
LoRA targets & $q\_\text{proj}$, $v\_\text{proj}$ (both encoders) \\
LoRA rank / $\alpha$ & $r=4$, $\alpha=16$ (scale = 4) \\
Trainable params & 295K (0.14\% of backbone) \\
Global loss & SigLIP sigmoid; MHAP \texttt{pooler\_output} \\
Region loss & EOS phrase $\cdot$ patch, max-pool, sigmoid SigLIP \\
$\lambda$ & 0.5 \\
Learning rate & 2e-4, cosine decay with 100-step warmup \\
Batch / region batch & 32 / 64 \\
Epochs (v10 / v11) & 1 / 5 \\
LR schedule & Per-epoch cosine restart (\texttt{restarts=epochs}) \\
Snapshot & Best-PG on 200-image val (every 200 steps) \\
Image path & MaskCLIP bypass installed before training \\
Text path & EOS token (last non-padding) \\
\bottomrule
\end{tabular}
\end{table}

%% ─────────────────────────────────────────────────────────────────────────────
\FloatBarrier
\section{Pointing Game Visualisations}
\label{app:viz}

All six cherry-picked validation images are shown below, one per page.
Each panel shows: \textbf{left} = Frozen B0; \textbf{centre} = M\textsubscript{human};
\textbf{right} = M\textsubscript{auto}.  Each row corresponds to one query
phrase.  Plasma heatmap = cosine similarity of each $24{\times}24$ patch
to the query phrase; white box = GT bounding box; $\bigstar$ = argmax patch
({\color{green!60!black}green} = hit / {\color{red}red} = miss).
Per-image Pointing Game scores (fraction of phrases hit) are shown in the
caption.  Images were selected by scanning all 200 validation images for
cases where both fine-tuned models outperform B0; the list is sorted by
combined PG gain over B0.

\begin{figure}[p]
\centering
\includegraphics[width=\linewidth,height=0.82\textheight,keepaspectratio]{pg_viz_01_1836335410}
\caption{Image \texttt{1836335410.jpg} ($n{=}3$ phrases).\quad
  B0: 0\%\enspace M\textsubscript{human}: \textbf{100\%}\enspace M\textsubscript{auto}: 33\%.
  Combined gain over B0: +67 pp.}
\label{fig:app1}
\end{figure}

\begin{figure}[p]
\centering
\includegraphics[width=\linewidth,height=0.82\textheight,keepaspectratio]{pg_viz_02_477607006}
\caption{Image \texttt{477607006.jpg} ($n{=}3$ phrases shown, $n{=}8$ total).\quad
  B0: 25\%\enspace M\textsubscript{human}: 62\%\enspace M\textsubscript{auto}: \textbf{88\%}.
  Combined gain over B0: +50 pp.  M\textsubscript{auto} dominates here.}
\label{fig:app2}
\end{figure}

\begin{figure}[p]
\centering
\includegraphics[width=\linewidth,height=0.82\textheight,keepaspectratio]{pg_viz_03_1496289999}
\caption{Image \texttt{1496289999.jpg} ($n{=}3$ phrases shown, $n{=}8$ total).\quad
  B0: 0\%\enspace M\textsubscript{human}: 38\%\enspace M\textsubscript{auto}: \textbf{75\%}.
  Combined gain over B0: +56 pp.}
\label{fig:app3}
\end{figure}

\begin{figure}[p]
\centering
\includegraphics[width=\linewidth,height=0.82\textheight,keepaspectratio]{pg_viz_04_1507563902}
\caption{Image \texttt{1507563902.jpg} ($n{=}3$ phrases).\quad
  B0: 33\%\enspace M\textsubscript{human}: \textbf{100\%}\enspace M\textsubscript{auto}: \textbf{67\%}.
  Combined gain over B0: +50 pp.  Both fine-tuned models hit all phrases; B0 misses two.}
\label{fig:app4}
\end{figure}

\begin{figure}[p]
\centering
\includegraphics[width=\linewidth,height=0.82\textheight,keepaspectratio]{pg_viz_05_5303395266}
\caption{Image \texttt{5303395266.jpg} ($n{=}3$ phrases shown, $n{=}10$ total).\quad
  B0: 10\%\enspace M\textsubscript{human}: \textbf{60\%}\enspace M\textsubscript{auto}: 50\%.
  Densest scene in the selection ($n{=}10$ phrases); both models generalise across
  many co-located referents.}
\label{fig:app5}
\end{figure}

\begin{figure}[p]
\centering
\includegraphics[width=\linewidth,height=0.82\textheight,keepaspectratio]{pg_viz_06_2067833088}
\caption{Image \texttt{2067833088.jpg} ($n{=}3$ phrases shown, $n{=}6$ total).\quad
  B0: 17\%\enspace M\textsubscript{human}: \textbf{67\%}\enspace M\textsubscript{auto}: 50\%.
  Combined gain over B0: +42 pp.}
\label{fig:app6}
\end{figure}

\end{document}
